% Source: https://artofproblemsolving.com/community/c3h2636758_intermediate_problems
% Stable identity: aops-topic-2636758-q14
\emph{One day.} The a.m.\ hours are $12,1,2,\dots,11$, with sum $1+2+\cdots+12=78$. The p.m.\ hours are the same, and the sum of their squares is $1^2+2^2+\cdots+12^2=\dfrac{12\cdot13\cdot25}{6}=650$. So on a day with number $d$ the clock rings $d(78+650)=728d$ times.

\emph{The year.} $2024$ is a leap year with $366=52\cdot7+2$ days. Starting on Monday 1 January, the first $52\cdot7=364$ days are $52$ full weeks, with day numbers adding up to $52(1+2+\cdots+7)=1456$. The last two days, 30 and 31 December, are a Monday and a Tuesday, adding $1+2=3$. So the day numbers of 2024 add up to $1459$, and the clock rings
\[728\cdot1459=\boxed{1062152}\]
times.
